\documentclass[10pt,a4paper]{amsart}
\usepackage[margin=19mm]{geometry}
\usepackage{amsmath,amssymb,mathtools}
\usepackage[scr=rsfs,cal=euler]{mathalfa}
\usepackage{xfrac}
\usepackage[unicode,hidelinks,bookmarks=false]{hyperref}
\newcommand{\ie}{i.e.\ }
\newcommand{\cf}{cf.\ }
\newcommand{\resp}{respectively\ }
\newcommand{\Subsection}{\subsection}
\providecommand{\nobreakdash}{\nobreak\hbox{-}}
\setlength{\emergencystretch}{2em}
\multlinegap0pt
\usepackage{amssymb}
\usepackage{xcolor}
\usepackage[normalem]{ulem}
\usepackage{tikz}
\usepackage{float}
\newtheorem{proposition}{Proposition}
\newtheorem{corollary}{Corollary}
\newtheorem{lemma}{Lemma}
\title{Expanding generalized fine rings}
\author{Ahmad Moussavi, Peter Danchev, Arash Javan,, Omid Hasanzadeh}
\date{Source: arXiv:2607.18125v1 (CC BY 4.0)}
\begin{document}
\maketitle

\noindent\emph{Source and licence.} Expanding generalized fine rings. Version arXiv:2607.18125v1 (CC BY 4.0), distributed by arXiv under the Creative Commons Attribution 4.0 International licence, http://creativecommons.org/licenses/by/4.0/, as recorded in the arXiv licence metadata for this exact version and shown on the abstract page https://arxiv.org/abs/2607.18125v1. Full licence notice: \texttt{licenses/arxiv-2607.18125-CC-BY-4.0.txt}.\par\smallskip

\noindent\textit{Editorial scope.} Counted unit: the article's lemma 2.4 (source0.tex lines 666--668). The article's complete original proof of it is delivered with it (source0.tex lines 670--758, one \texttt{proof} environment of 89 lines). No mathematics is written, repaired or re-derived: the statement and the proof are the article's own lines, in the article's own notation, and no retained line is substituted or re-flowed.  Retained uncounted as the source-local context the proof needs: lines 269--277; lines 660--662.  Nothing was omitted from any retained span, so the package carries no omission record.  No retained line contains a citation, so the wrapper carries no bibliography and no bibliography entry.  No retained line contains a figure environment, an image-inclusion command or drawing code, so no image is delivered and the package declares no asset dependency.  Editorial formatting only, in addition: an AMS \texttt{amsart} preamble that restates the article's own declarations and macros used by the retained lines, namely the environments \texttt{proposition}, \texttt{corollary}, \texttt{lemma} on separate counters, together with the standard packages those lines need.  The retained environments print in this document as Proposition 1, Corollary 1, Lemma 1 in order of appearance, whereas the article prints the same items with the numbering of its own full document.  The complete native source of the article as distributed in the arXiv e-print for this exact version is bundled unchanged as \texttt{sources/205628/source0.tex}.
\begin{proposition}\label{2}
For any ring \(R\), the following three points hold:

(1) If \(R \neq 0\), then \(\sqrt{J}F(R) \cap C(R) = U(C(R))\).

(2) \(v \in U(R)\) implies \(v\sqrt{J}F(R)v^{-1} = \sqrt{J}F(R)\). If, moreover, \(v \in U(R) \cap C(R)\), then \(v\sqrt{J}F(R) = \sqrt{J}F(R)\).

(3) If \(R \neq 0\) and \(\sqrt{J(R)} \subseteq J(R)\) (i.e., \(\sqrt{J(R)} = J(R)\)), then \(\sqrt{J}F(R) = U(R)\).
\end{proposition}

\begin{corollary}\label{Block}
Let \( M = (A_{ij}) \in R = M_n(S) \) be a block matrix with each diagonal block \( A_{ii} \) \(\sqrt{J}\)-fine. Then, \( M \) belongs to \(\sqrt{J}F(R)\).
\end{corollary}

\begin{lemma}\label{2.4}
If \( R \) is a generalized $\sqrt{J}$-fine ring, then \( M_2(R) \) is also a generalized $\sqrt{J}$-fine ring.
\end{lemma}

\begin{proof}
Let
	\[
	A := \begin{pmatrix} a & b \\ c & d \end{pmatrix} \in M_2(R) \setminus M_2(J(R)).
	\]
We consider four possible cases based on how many entries of \( A \) are outside \( J(R) \).

\medskip
	
\textbf{Case 1: All four entries are not in \( J(R) \).} Then, both \( a \) and \( d \) are $\sqrt{J}$-fine in \( R \). By Corollary \ref{Block}, \( A \) is $\sqrt{J}$-fine.

\medskip
	
\textbf{Case 2: Exactly three entries are not in \( J(R) \).} So, exactly one entry is in \( J(R) \).

\medskip
	
If \( a \in J(R) \), then \( A \) is $\sqrt{J}$-fine if and only if \( \begin{pmatrix} 0 & b \\ c & d \end{pmatrix} \) is $\sqrt{J}$-fine. So, we may assume \( a = 0 \).
Write
	\[
	A_1 = \begin{pmatrix} 0 & 1 \\ 1 & 1 \end{pmatrix} \begin{pmatrix} 0 & b \\ c & d \end{pmatrix} \begin{pmatrix} -1 & 1 \\ 1 & 0 \end{pmatrix}
	= \begin{pmatrix} -c+d & c \\ b-c+d & c \end{pmatrix}.
	\]
If \( -c+d \not\in J(R) \), then \( A_1 \) and hence \( A \) is $\sqrt{J}$-fine by virtue of Corollary \ref{Block} and Proposition \ref{2}(2). So, assume \( -c+d \in J(R) \). Then, \( A \) is $\sqrt{J}$-fine if and only if \( A_1 \) is $\sqrt{J}$-fine, which amounts to \( \begin{pmatrix} 0 & c \\ b & c \end{pmatrix} \) being $\sqrt{J}$-fine. So, we may assume
$
	A_1 = \begin{pmatrix} 0 & c \\ b & c \end{pmatrix}.
$
Now, let
	\[
	A_2 := \begin{pmatrix} 1 & 1 \\ 0 & 1 \end{pmatrix} \begin{pmatrix} 0 & c \\ b & c \end{pmatrix} \begin{pmatrix} 1 & -1 \\ 0 & 1 \end{pmatrix}
	= \begin{pmatrix} b & -b+2c \\ b & -b+c \end{pmatrix}.
	\]
If \( -b+c \not\in J(R) \), then \( A_2 \) and hence \( A_1 \) is $\sqrt{J}$-fine by usage of Corollary \ref{Block} and Proposition \ref{2}(2). So, assume \( -b+c \in J(R) \), i.e., \( b = c + j \) for some \( j \in J(R) \). Thus, \( A_1 \) is $\sqrt{J}$-fine if and only if \( \begin{pmatrix} 0 & c \\ c & c \end{pmatrix} \) is $\sqrt{J}$-fine. Since \( c \not\in J(R) \), we can write \( c = u + x \) with \( u \in U(R) \) and \( x \in \sqrt{J(R)} \). Therefore,
	\[
	\begin{pmatrix} 0 & c \\ c & c \end{pmatrix}
	= \begin{pmatrix} 0 & u \\ u & c \end{pmatrix} + \begin{pmatrix} 0 & x \\ x & 0 \end{pmatrix},
	\]
which is a unit plus an element from $\sqrt{J}$. So, \( A \) is $\sqrt{J}$-fine.
	
If \( b \in J(R) \), then \( a \) and \( d \) are $\sqrt{J}$-fine in \( R \), so \( A \) is $\sqrt{J}$-fine by Corollary \ref{Block}.
	
The cases \( c \in J(R) \) or \( d \in J(R) \) are quite similar, so we skip their evidences.

\medskip
	
\textbf{Case 3: Exactly two entries are not in \( J(R) \).}

\medskip
	
If \( a \not\in J(R) \) and \( d \not\in J(R) \), then both \( a \) and \( d \) are $\sqrt{J}$-fine, so \( A \) is $\sqrt{J}$-fine by Corollary \ref{Block}.
	
If \( a \not\in J(R) \) and \( b \not\in J(R) \), then \( b \) is $\sqrt{J}$-fine in \( R \), and \( A \) is $\sqrt{J}$-fine if and only if \( \begin{pmatrix} a & b \\ 0 & 0 \end{pmatrix} \) is $\sqrt{J}$-fine. Write \( b = u + x \) with \( u \in U(R) \) and \( x \in \sqrt{J(R)} \). Thus,
	\[
	\begin{pmatrix} a & b \\ 0 & 0 \end{pmatrix}
	= \begin{pmatrix} a & u \\ -1 & 0 \end{pmatrix} + \begin{pmatrix} 0 & x \\ 1 & 0 \end{pmatrix},
	\]
is a sum of a unit and an element from $\sqrt{J}$. So, \( A \) is $\sqrt{J}$-fine.
	
If \( b \not\in J(R) \) and \( c \not\in J(R) \), then \( A \) is $\sqrt{J}$-fine if and only if \( \begin{pmatrix} 0 & b \\ c & 0 \end{pmatrix} \) is $\sqrt{J}$-fine. But,
	\[
	\begin{pmatrix} 0 & 1 \\ 1 & 1 \end{pmatrix} \begin{pmatrix} 0 & b \\ c & 0 \end{pmatrix} \begin{pmatrix} -1 & 1 \\ 1 & 0 \end{pmatrix}
	= \begin{pmatrix} -c & c \\ b-c & c \end{pmatrix},
	\]
which is $\sqrt{J}$-fine by Corollary \ref{Block}. Hence, \( A \) is $\sqrt{J}$-fine with Proposition \ref{2}(2) at hand.
	
The remaining subcases (\( a,c \not\in J(R) \), or \( b,d \not\in J(R) \), or \( c,d \not\in J(R) \)) are rather analogous, and thus we omit their verifications.

\medskip
	
\textbf{Case 4: Exactly one entry is not in \( J(R) \).}

\medskip
	
If \( a \not\in J(R) \), then \( A \) is $\sqrt{J}$-fine if and only if \( \begin{pmatrix} a & 0 \\ 0 & 0 \end{pmatrix} \) is $\sqrt{J}$-fine. But,
	\[
	\begin{pmatrix} 1 & 0 \\ 1 & 1 \end{pmatrix} \begin{pmatrix} a & 0 \\ 0 & 0 \end{pmatrix} \begin{pmatrix} 1 & 0 \\ -1 & 1 \end{pmatrix}
	= \begin{pmatrix} a & 0 \\ a & 0 \end{pmatrix},
	\]
which is $\sqrt{J}$-fine with the help of Case 2. So, \( A \) is $\sqrt{J}$-fine.
	
If \( b \not\in J(R) \), then \( b \) is $\sqrt{J}$-fine, so write \( b = u + x \) with \( u \in U(R) \) and \( x \in \sqrt{J(R)} \). Then, \( A \) is $\sqrt{J}$-fine if and only if \( \begin{pmatrix} 0 & b \\ 0 & 0 \end{pmatrix} \) is $\sqrt{J}$-fine. But,
	\[
	\begin{pmatrix} 0 & b \\ 0 & 0 \end{pmatrix}
	= \begin{pmatrix} 0 & u \\ -1 & 0 \end{pmatrix} + \begin{pmatrix} 0 & x \\ 1 & 0 \end{pmatrix},
	\]
is a sum of a unit and an element from $\sqrt{J}$. So, \( A \) is $\sqrt{J}$-fine.
	
Finally, in all cases, \( A \) is $\sqrt{J}$-fine, and consequently \( M_2(R) \) is generalized $\sqrt{J}$-fine.
\end{proof}

\end{document}
